% short/03_processing.tex: from the raw signals to a standard movement, in the order a reader needs: the even 60 Hz
% clock (native rates: nothing lost), smoothing (zero-phase Butterworth, gradual), up and the turn about the vertical
% (gravity removed, the gravity frame and the yaw), the iPhone accelerometer sign, the body frame (Figure 5, TikZ:
% two horizontal directions set by where the screen faces, so they turn with the phone's heading but not with the
% grip), cutting out the movement and time-normalising it to 64 points. Numbers and settings only through \R keys.
\section{From raw signals to a standard movement}\label{sec:processing}

\paragraph{An even clock.} A web page receives the motion samples about \R{describe/rate_hz_median_whole} times per second:
per session (one phone on one visit), the median rate ranged from \R{describe/rate_hz_min}~Hz (on
\R{describe/rate_hz_min_platform}) to \R{describe/rate_hz_max}~Hz (on \R{describe/rate_hz_max_platform}). The samples do not arrive at perfectly even times, and some iPhone recordings have gaps of about a third of a
second (\R{describe/ios_long_gap_attempts} recordings, median \R{describe/ios_long_gap_median_ms}~ms; the lowest average
rate of a single recording was \R{describe/rate_hz_recording_min}~Hz), so we put every recording on an even
\R{describe/fs_hz}~Hz clock by linear interpolation, bridging gaps of up to \R{describe/max_gap_ms}~ms (a trial with a
longer gap near the movement was not kept, Section~\ref{sec:data}). This loses
essentially nothing: the phones delivered about as many samples as the clock has, and the arm movement is slow
(\R{describe/gyro_power_below_5hz} of the gyroscope's signal power lies below 5~Hz).

\paragraph{Smoothing.} We smooth every channel with a standard low-pass filter (a Butterworth filter of order
\R{describe/lowpass_order}, cut-off \R{describe/lowpass_hz}~Hz), run once forward and once backward in time so that it
does not shift peaks in time. It is not a sharp cut: it damps fast wiggles gradually, keeping
\R{describe/lowpass_gain_10hz} of the amplitude of a 10~Hz oscillation, \R{describe/lowpass_gain_15hz} at
\R{describe/lowpass_hz}~Hz and \R{describe/lowpass_gain_20hz} at 20~Hz. During the movement,
\R{describe/gyro_power_below_lowpass} of the gyroscope's power and \R{describe/acc_power_below_lowpass} of the
accelerometer's lie below \R{describe/lowpass_hz}~Hz, so the filter removes little of the movement itself.

\paragraph{Up, and the turn about the vertical.} The phone's own axes turn with the grip, so we express the motion
relative to the vertical. When the phone is still, the accelerometer feels only gravity, which shows which way is down;
we measure it in the still \R{describe/rest_ms}~ms just before the movement (at arm's length) and just after it (at the
nose). During the movement the accelerometer also feels the movement itself, so there we carry `down' with the
gyroscope: it tells how far the phone has rotated since the still moment, and we rotate `down' back by the same amount,
sample by sample, forward from the start and backward from the end, and blend the two. (We did not use the phone's
compass, the magnetometer, so we know where up is but not which way the person faces; phones also report their own
estimate of gravity, but we computed ours.) Carried through the whole movement, the `down' of the start arrives within
a median of \R{describe/closure_median_deg}$^\circ$ of the `down' measured at the nose.

With up known at every sample, we subtract gravity from the accelerometer, so that every acceleration from here
on is that of the movement alone (zero when the phone is still), and we split each rotation and acceleration into a vertical and a horizontal part. This
gives \R{final/cnn_n_gravity} \emph{gravity-frame} channels (a frame is a set of directions to measure along; this one
is fixed by gravity alone) that do not depend on how the phone sits in the hand: the total rotation speed, the total
acceleration, the vertical and the horizontal acceleration, the tilting speed (how fast the phone tilts about any
horizontal axis, a size without a direction), the yaw rate, the yaw and the tilt since the start (how far the phone
has tipped relative to the vertical since the start of the movement). The \emph{yaw rate} is how fast
the phone turns about the vertical, like a door on its hinges, counted positive counter-clockwise seen from above
(Figure~\ref{fig:frames}c); the \emph{yaw} adds it up from the start of the movement, that is, how far the phone has
turned about the vertical so far.

\begin{figure}[tp]
\centering
% short/diagram_frames.tex: the phone's own axes against the directions the analysis measures in (short.tex,
% Section 3). TikZ at a fixed size (about 14 cm wide, explicit 8/9-pt fonts), used at full text width. Seen from the
% person, who looks at the screen: "towards the face" points out of the page (a circled dot).
% (a) upright grip: x to the right, y up. (b) sideways grip with the top of the phone to the person's LEFT: y to the
% left, x up (x = y cross z, with z towards the face). The body-frame directions up / right / towards the face are the
% same in (a) and (b). (c) seen from above: the turn about the vertical (yaw rate), positive counter-clockwise.
\begin{tikzpicture}[
    font=\fontsize{8}{9.5}\selectfont,
    phone/.style={draw=black!70, fill=black!4, rounded corners=3pt, line width=0.7pt},
    dev/.style={-{Latex[length=2mm,width=1.7mm]}, black!45, line width=1.1pt},
    body/.style={-{Latex[length=2mm,width=1.7mm]}, black, line width=1.1pt},
    devlab/.style={text=black!55, inner sep=1.5pt, font=\fontsize{9}{10}\selectfont},
    panel/.style={anchor=north, font=\fontsize{8.5}{10}\selectfont, align=center}]
  % \outofpage{point}{colour}: an arrow pointing out of the page, towards the viewer (a circled dot)
  \def\outofpage#1#2{\fill[white] #1 circle (1.5mm); \draw[#2, line width=0.8pt] #1 circle (1.5mm);
    \fill[#2] #1 circle (0.45mm);}
  % ---------------------------------------------------------------- (a) upright
  \begin{scope}
    \draw[phone] (-0.55,-1.0) rectangle (0.55,1.0);
    \draw[dev] (0,0) -- (1.0,0) node[devlab, right] {$x$};
    \draw[dev] (0,0) -- (0,1.45) node[devlab, above] {$y$};
    \outofpage{(0,0)}{black!45}
    \node[devlab, anchor=north west] at (0.08,-0.10) {$z$};
    \begin{scope}[shift={(2.15,-0.75)}]
      \draw[body] (0,0) -- (0,1.15) node[above, inner sep=1.5pt] {up};
      \draw[body] (0,0) -- (1.15,0) node[right, inner sep=1.5pt] {right};
      \outofpage{(0,0)}{black}
      \node[anchor=north, inner sep=2pt] at (0,-0.16) {towards the face};
    \end{scope}
    \node[panel] at (1.35,-1.55) {(a) phone upright};
  \end{scope}
  % ---------------------------------------------------------------- (b) sideways, top of the phone to the LEFT
  \begin{scope}[shift={(6.15,0)}]
    \draw[phone] (-1.0,-0.55) rectangle (1.0,0.55);
    \draw[dev] (0,0) -- (0,1.45) node[devlab, above] {$x$};
    \draw[dev] (0,0) -- (-1.45,0) node[devlab, left] {$y$};
    \outofpage{(0,0)}{black!45}
    \node[devlab, anchor=north west] at (0.08,-0.10) {$z$};
    \begin{scope}[shift={(1.75,-0.75)}]
      \draw[body] (0,0) -- (0,1.15) node[above, inner sep=1.5pt] {up};
      \draw[body] (0,0) -- (1.15,0) node[right, inner sep=1.5pt] {right};
      \outofpage{(0,0)}{black}
      \node[anchor=north, inner sep=2pt] at (0,-0.16) {towards the face};
    \end{scope}
    \node[panel] at (0.55,-1.55) {(b) phone sideways (top to the left)};
  \end{scope}
  % ---------------------------------------------------------------- (c) seen from above: the turn about the vertical
  \begin{scope}[shift={(12.0,0.05)}]
    \draw[black!20, line width=0.6pt] (0,0) circle (0.85);
    \draw[body] (-30:0.85) arc (-30:240:0.85);
    \outofpage{(0,0)}{black}
    \node[anchor=north, inner sep=2pt] at (0,-0.18) {up};
    \node[font=\fontsize{11}{12}\selectfont\bfseries] at (1.0,0.75) {$+$};
    \node[panel] at (0,-1.55) {(c) seen from above};
  \end{scope}
\end{tikzpicture}%

\caption{The phone's own axes (grey) and the directions we measure in (black; this section explains them), as the
person sees the phone, so that $z$ and `towards the face' point out of the page. (a, b) Turning the phone in the hand
from upright to sideways swaps its $x$ and $y$ axes, but up, right and towards the face stay the same. (c) Seen from
above, up points out of the page: the yaw rate, the turn about the vertical, is positive counter-clockwise.}
\label{fig:frames}
\end{figure}

\paragraph{The iPhone sign.} Measuring `down' needs the right sign. iPhones report the accelerometer with the opposite
sign to Android phones on all three axes (the gyroscope has the same sign on both), so on an iPhone `down' would come
out as `up', every turn about the vertical would be reversed and a left hand would look like a right hand. We therefore
multiplied every iPhone accelerometer value by $-1$ and checked each session on its own recordings: while the phone is
held still before a movement, the corrected reading must point up along the axis that the grip puts on top. All
\R{describe/accel_sign_flipped_ios} iPhone sessions needed the flip, and none of the \R{describe/sessions_android}
Android sessions did.

\paragraph{The body frame.} The gravity frame has no horizontal direction, so it cannot tell which way the phone moves
or tilts sideways; the phone's own axes can, but they change with the grip. Four more channels therefore use two
horizontal directions set by where the screen faces (Figure~\ref{fig:frames}): `towards the face', the direction the
screen faces, made horizontal, and `right', horizontal and perpendicular to it, towards the person's right. Since the
screen faces the person, these are in practice the directions of the person's body (hence the name). They turn with the phone when it turns about the vertical, but not when it is turned
within the plane of its screen, which is all that a change of grip does: held upright or sideways, the screen still
faces the person, so up, right and towards the face, and every channel measured along them, stay the same, while the
phone's own $x$ and $y$ axes swap places (Figure~\ref{fig:frames}a, b). Along these directions we measure the roll
rate (rotation about the towards-the-face direction) and the pitch rate (rotation about the right direction), which
split the tilting speed into its two horizontal directions, with signs, and the sideways acceleration and the
acceleration towards the face. The final model reads these \R{final/cnn_n_body}
\emph{body-frame} channels together with the \R{final/cnn_n_gravity} gravity-frame channels.

\paragraph{Cutting out the movement, and time-normalising it.} The models read only the movement to the nose, between
the app's two markers. The start marker is the last still moment before the phone's rotation or acceleration rose
clearly above its resting level, a median of \R{describe/reaction_median_ms}~ms after the start beep; the end marker is
the moment the phone came to rest at the nose. A median movement lasts \R{describe/duration_median_s}~s, that is,
\R{describe/approach_samples_median} samples at \R{describe/fs_hz}~Hz. Between the two markers we compute each channel
at \R{describe/n_norm} points evenly spaced in time, by interpolating between the samples, whatever the duration: the
first point is the start and the last is the arrival at the nose. Time then runs as a percentage of the movement,
from 0 to 100\,\%, so movements of different speeds line up point by point, and every movement becomes an input of
the same size.
